\documentclass[11pt]{article}

\usepackage[letterpaper,margin=1in]{geometry}
\usepackage[T1]{fontenc}
\usepackage{lmodern}
\usepackage{microtype}
\usepackage{amsmath,amssymb,mathtools}
\usepackage{booktabs}
\usepackage[hidelinks]{hyperref}

\title{Section 8 and Predator 8.002:\\
From Counting Bounds to ML-Guided Proof Search}
\author{Brian Tenneson}
\date{August 1, 2026}

\newcommand{\prcom}{\texttt{prcom}}

\begin{document}

\maketitle

\section{The connection}

Section 8 of \emph{Depths of a Simulation}, Version 45, bounds the
number of possible length-$n$ proof sequences.  In its notation,

\begin{equation}
P_F(\Gamma,\varphi,n)
\leq
\prod_{j=1}^{n}
\left(
|\Gamma|+\sum_{k=1}^{a}\kappa_k(j-1)^k
\right).
\label{eq:section8-bound}
\end{equation}

Here $\kappa_k$ is the number of $k$-ary rules, $a$ is the maximum
rule arity, and $P_F(\Gamma,\varphi,n)$ is the number of proofs of
$\varphi$ from $\Gamma$ having length exactly $n$.  The theorem is
deliberately a crude cardinal upper bound: it counts many locally
describable line-construction sequences that are not genuine proofs,
as well as redundant constructions and constructions irrelevant to
the target.

Predator 8.002 provides an experimental realization of the structural
observation following this theorem.  Proofs are not structureless
passwords.  They contain reusable lemmas, legal intermediate states,
search-order information, and branching structure.  Machine learning
can exploit that structure without changing the formal rules or the
set of legal proofs.

\section{A schema-level illustration for \prcom}

In the frozen Metamath environment used by Predator, the portion of
the database strictly before \prcom{} contains the following data:

\begin{center}
\begin{tabular}{@{}lr@{}}
\toprule
Quantity & Value \\
\midrule
Earlier assertion schemas & 4,696 \\
Maximum raw mandatory-hypothesis arity & 41 \\
Floating hypotheses of \prcom & 2 \\
Labels in the generated Metamath certificate & 28 \\
Principal logical assertion applications & 4 \\
\bottomrule
\end{tabular}
\end{center}

The four principal logical applications in the certificate are

\[
\texttt{uncom},\qquad
\texttt{df-pr},\qquad
\texttt{df-pr},\qquad
\texttt{3eqtr4i}.
\]

Under a finite schema-level interpretation of
Equation~\eqref{eq:section8-bound}, with the 41 zero-arity schemas
absorbed into the initial resources, the crude 28-line upper bound is
approximately

\begin{equation}
P_F(\Gamma,\prcom,28)
\lesssim 1.69\times 10^{1167}.
\label{eq:raw-proxy}
\end{equation}

This number is not an estimate of the actual number of proofs of
\prcom, and it is not the number of states that Predator must search.
It illustrates the scale of the combinatorial possibility space
described by the Section 8 bound.

At the logical-only level, there are 4,643 earlier logical assertion
schemas.  Treating the 2,157 assertions with no essential hypotheses
as base options produces an analogous four-logical-application proxy
of approximately

\begin{equation}
5.1\times 10^{19}.
\label{eq:logical-proxy}
\end{equation}

Equation~\eqref{eq:logical-proxy} is again an illustrative schema-level
bound, not an exact target-specific proof count.  Metamath assertions
are schematic and admit substitutions, while Equation
\eqref{eq:section8-bound} is stated for finitary partial-function
rules.  A completely literal translation would therefore need to
make the treatment of substitutions and instantiated rules explicit.

\section{From counting to learned search measures}

Section 8 concerns how many proof sequences may exist.  Predator
places a learned, nonuniform ordering over legal proof-search moves.
Conceptually, if $s_\theta(a\mid S)$ is the learned score assigned to
a legal assertion application $a$ in proof state $S$, a temperature
parameter $T$ induces a distribution of the form

\begin{equation}
p_{\theta,T}(a\mid S)
=
\frac{\exp\!\left(s_\theta(a\mid S)/T\right)}
{\displaystyle\sum_{b\in A(S)}
 \exp\!\left(s_\theta(b\mid S)/T\right)},
\label{eq:policy}
\end{equation}

where $A(S)$ is the set of legal applications in $S$.  This is an
idealized description of the implemented combination of learned
scores, seeded stochastic ordering, novelty, rarity, and
counterfactual exploration.  A complete proof path
$\pi=(a_1,\ldots,a_n)$ then receives an idealized path probability

\begin{equation}
p_{\theta,T}(\pi)
=
\prod_{j=1}^{n}p_{\theta,T}(a_j\mid S_j).
\label{eq:path-probability}
\end{equation}

Thus Section 8 supplies a counting measure on the space of possible
proof constructions, while Predator supplies a learned search measure
on the legal part of that space.  Machine learning does not remove
legal proofs or create new inference rules.  It reallocates finite
search attention.

The creativity controls regulate how concentrated this attention is.
Low temperature concentrates probability around highly ranked moves;
higher temperature and counterfactual selection preserve some search
mass for lower-ranked legal alternatives.  The Metamath verifier is
independent of both mechanisms and remains the final authority on
whether the selected path is a proof.

\section{The controlled \prcom{} result}

After parsing and unification, Predator found 150 legal assertion
applications at the root of \prcom.  The useful equality-chaining
assertion \texttt{3eqtr4i} behaved as follows:

\begin{center}
\begin{tabular}{@{}lr@{}}
\toprule
Measurement & Result \\
\midrule
Unguided diagnostic rank of \texttt{3eqtr4i} & 134 of 150 \\
ML-guided raw rank of \texttt{3eqtr4i} & 64 of 150 \\
Approximate ordinary explorer capacity & 60 \\
Guided first verified proof & 2,633 expansions \\
Unguided result at the common cutoff & no proof in 5,000 expansions \\
\bottomrule
\end{tabular}
\end{center}

The ML score moved \texttt{3eqtr4i} close to the ordinary explorer
cutoff.  Counterfactual creativity then admitted the slightly
out-of-cap candidate, after which the formal search completed all
subgoals and emitted a certificate.

Let $E_M(\varphi)$ denote the number of search-state expansions before
prover $M$ first discovers a verified proof of $\varphi$, with
$E_M(\varphi)=\infty$ if it never does.  The experiment established

\[
E_{\mathrm{ML}}(\prcom)=2633
\qquad\text{and}\qquad
E_{\mathrm{unguided}}(\prcom)>5000.
\]

Consequently, the censored fixed-instance search-effort ratio satisfies

\begin{equation}
\frac{E_{\mathrm{unguided}}(\prcom)}
     {E_{\mathrm{ML}}(\prcom)}
>
\frac{5000}{2633}
\approx 1.90.
\label{eq:fixed-separation}
\end{equation}

The true ratio is unknown because the unguided condition did not find
a proof before its stopping rule.

The expansion count in Equation~\eqref{eq:fixed-separation} is not the
same quantity as the proof length $n$ in Equation
\eqref{eq:section8-bound}.  Proof length measures the size of a
certificate or derivation under a declared convention.  Predator's
budget counts search-state expansions performed while trying to find
such a derivation.  The distinction is essential: ML reduced discovery
effort, not the formal length or logical validity conditions of the
certificate.

\section{Interpretation}

The relationship can be summarized as follows:

\begin{center}
\begin{tabular}{@{}p{0.28\textwidth}p{0.62\textwidth}@{}}
\toprule
Component & Function \\
\midrule
Section 8 bound & Bounds the cardinal size of proof-construction space \\
Parsing and unification & Remove syntactically or formally illegal applications \\
ML policy & Prioritizes structurally promising legal applications \\
Creativity controls & Preserve alternatives when learned ranking is imperfect \\
Metamath verifier & Certifies the complete selected derivation \\
\bottomrule
\end{tabular}
\end{center}

Predator 8.002 did not lower the Section 8 upper bound and did not
change which proofs were legal.  It concentrated a finite expansion
budget on a smaller and more productive portion of the legal search
space.  The experiment is therefore a concrete finite-instance
demonstration of Section 8's concluding principle: search procedures
can be compared mathematically and experimentally by how quickly they
reach certified solutions.

It is not yet an asymptotic speedup theorem.  Such a theorem would
require a family of targets, a scaling parameter, and proved upper and
lower bounds for the competing procedures.  The \prcom{} result is a
controlled proof of concept showing that the proposed mechanism can
produce a genuine search-effort separation while formal soundness
remains unchanged.

For the HaloProof benchmark, the same framework suggests the relevant
question.  Its formal proof space may be enormous, but learned guidance
could concentrate search mass around productive constructions, while
controlled creativity prevents an imperfect model from assigning
effectively zero attention to the necessary route.  The benchmark will
test whether this finite-instance mechanism transfers to a longer and
substantially more difficult derivation.

\end{document}
